

\documentclass[journal]{Imperial_lab_report}

% *** MISC UTILITY PACKAGES ***
%
%\usepackage{ifpdf}
% Heiko Oberdiek's ifpdf.sty is very useful if you need conditional
% compilation based on whether the output is pdf or dvi.
% usage:
% \ifpdf
%   % pdf code
% \else
%   % dvi code
% \fi
% The latest version of ifpdf.sty can be obtained from:
% http://www.ctan.org/pkg/ifpdf
% Also, note that IEEEtran.cls V1.7 and later provides a builtin
% \ifCLASSINFOpdf conditional that works the same way.
% When switching from latex to pdflatex and vice-versa, the compiler may
% have to be run twice to clear warning/error messages.


% *** CITATION PACKAGES ***
%
%\usepackage{cite}
% cite.sty was written by Donald Arseneau
% V1.6 and later of IEEEtran pre-defines the format of the cite.sty package
% \cite{} output to follow that of the IEEE. Loading the cite package will
% result in citation numbers being automatically sorted and properly
% "compressed/ranged". e.g., [1], [9], [2], [7], [5], [6] without using
% cite.sty will become [1], [2], [5]--[7], [9] using cite.sty. cite.sty's
% \cite will automatically add leading space, if needed. Use cite.sty's
% noadjust option (cite.sty V3.8 and later) if you want to turn this off
% such as if a citation ever needs to be enclosed in parenthesis.
% cite.sty is already installed on most LaTeX systems. Be sure and use
% version 5.0 (2009-03-20) and later if using hyperref.sty.
% The latest version can be obtained at:
% http://www.ctan.org/pkg/cite
% The documentation is contained in the cite.sty file itself.


% *** GRAPHICS RELATED PACKAGES ***
%
\ifCLASSINFOpdf
\usepackage[pdftex]{graphicx}
  % declare the path(s) where your graphic files are
  % \graphicspath{{../pdf/}{../jpeg/}}
  % and their extensions so you won't have to specify these with
  % every instance of \includegraphics
  % \DeclareGraphicsExtensions{.pdf,.jpeg,.png}
\else
  % or other class option (dvipsone, dvipdf, if not using dvips). graphicx
  % will default to the driver specified in the system graphics.cfg if no
  % driver is specified.
  % \usepackage[dvips]{graphicx}
  % declare the path(s) where your graphic files are
  % \graphicspath{{../eps/}}
  % and their extensions so you won't have to specify these with
  % every instance of \includegraphics
  % \DeclareGraphicsExtensions{.eps}
\fi
% graphicx was written by David Carlisle and Sebastian Rahtz. It is
% required if you want graphics, photos, etc. graphicx.sty is already
% installed on most LaTeX systems. The latest version and documentation
% can be obtained at: 
% http://www.ctan.org/pkg/graphicx
% Another good source of documentation is "Using Imported Graphics in
% LaTeX2e" by Keith Reckdahl which can be found at:
% http://www.ctan.org/pkg/epslatex
%
% latex, and pdflatex in dvi mode, support graphics in encapsulated
% postscript (.eps) format. pdflatex in pdf mode supports graphics
% in .pdf, .jpeg, .png and .mps (metapost) formats. Users should ensure
% that all non-photo figures use a vector format (.eps, .pdf, .mps) and
% not a bitmapped formats (.jpeg, .png). The IEEE frowns on bitmapped formats
% which can result in "jaggedy"/blurry rendering of lines and letters as
% well as large increases in file sizes.
%
% You can find documentation about the pdfTeX application at:
% http://www.tug.org/applications/pdftex


% *** MATH PACKAGES ***
%
\usepackage{amsmath}
% A popular package from the American Mathematical Society that provides
% many useful and powerful commands for dealing with mathematics.
%
% Note that the amsmath package sets \interdisplaylinepenalty to 10000
% thus preventing page breaks from occurring within multiline equations. Use:
%\interdisplaylinepenalty=2500
% after loading amsmath to restore such page breaks as IEEEtran.cls normally
% does. amsmath.sty is already installed on most LaTeX systems. The latest
% version and documentation can be obtained at:
% http://www.ctan.org/pkg/amsmath





% *** SPECIALIZED LIST PACKAGES ***
%
%\usepackage{algorithmic}
% algorithmic.sty was written by Peter Williams and Rogerio Brito.
% This package provides an algorithmic environment fo describing algorithms.
% You can use the algorithmic environment in-text or within a figure
% environment to provide for a floating algorithm. Do NOT use the algorithm
% floating environment provided by algorithm.sty (by the same authors) or
% algorithm2e.sty (by Christophe Fiorio) as the IEEE does not use dedicated
% algorithm float types and packages that provide these will not provide
% correct IEEE style captions. The latest version and documentation of
% algorithmic.sty can be obtained at:
% http://www.ctan.org/pkg/algorithms
% Also of interest may be the (relatively newer and more customizable)
% algorithmicx.sty package by Szasz Janos:
% http://www.ctan.org/pkg/algorithmicx


% *** ALIGNMENT PACKAGES ***
%
%\usepackage{array}
% Frank Mittelbach's and David Carlisle's array.sty patches and improves
% the standard LaTeX2e array and tabular environments to provide better
% appearance and additional user controls. As the default LaTeX2e table
% generation code is lacking to the point of almost being broken with
% respect to the quality of the end results, all users are strongly
% advised to use an enhanced (at the very least that provided by array.sty)
% set of table tools. array.sty is already installed on most systems. The
% latest version and documentation can be obtained at:
% http://www.ctan.org/pkg/array


% IEEEtran contains the IEEEeqnarray family of commands that can be used to
% generate multiline equations as well as matrices, tables, etc., of high
% quality.


% *** SUBFIGURE PACKAGES ***
%\ifCLASSOPTIONcompsoc
%  \usepackage[caption=false,font=normalsize,labelfont=sf,textfont=sf]{subfig}
%\else
%  \usepackage[caption=false,font=footnotesize]{subfig}
%\fi
% subfig.sty, written by Steven Douglas Cochran, is the modern replacement
% for subfigure.sty, the latter of which is no longer maintained and is
% incompatible with some LaTeX packages including fixltx2e. However,
% subfig.sty requires and automatically loads Axel Sommerfeldt's caption.sty
% which will override IEEEtran.cls' handling of captions and this will result
% in non-IEEE style figure/table captions. To prevent this problem, be sure
% and invoke subfig.sty's "caption=false" package option (available since
% subfig.sty version 1.3, 2005/06/28) as this is will preserve IEEEtran.cls
% handling of captions.
% Note that the Computer Society format requires a larger sans serif font
% than the serif footnote size font used in traditional IEEE formatting
% and thus the need to invoke different subfig.sty package options depending
% on whether compsoc mode has been enabled.
%
% The latest version and documentation of subfig.sty can be obtained at:
% http://www.ctan.org/pkg/subfig


\usepackage{caption}

% *** FLOAT PACKAGES ***
%
%\usepackage{fixltx2e}
% fixltx2e, the successor to the earlier fix2col.sty, was written by
% Frank Mittelbach and David Carlisle. This package corrects a few problems
% in the LaTeX2e kernel, the most notable of which is that in current
% LaTeX2e releases, the ordering of single and double column floats is not
% guaranteed to be preserved. Thus, an unpatched LaTeX2e can allow a
% single column figure to be placed prior to an earlier double column
% figure.
% Be aware that LaTeX2e kernels dated 2015 and later have fixltx2e.sty's
% corrections already built into the system in which case a warning will
% be issued if an attempt is made to load fixltx2e.sty as it is no longer
% needed.
% The latest version and documentation can be found at:
% http://www.ctan.org/pkg/fixltx2e


%\usepackage{stfloats}
% stfloats.sty was written by Sigitas Tolusis. This package gives LaTeX2e
% the ability to do double column floats at the bottom of the page as well
% as the top. (e.g., "\begin{figure*}[!b]" is not normally possible in
% LaTeX2e). It also provides a command:
%\fnbelowfloat
% to enable the placement of footnotes below bottom floats (the standard
% LaTeX2e kernel puts them above bottom floats). This is an invasive package
% which rewrites many portions of the LaTeX2e float routines. It may not work
% with other packages that modify the LaTeX2e float routines. The latest
% version and documentation can be obtained at:
% http://www.ctan.org/pkg/stfloats
% Do not use the stfloats baselinefloat ability as the IEEE does not allow
% \baselineskip to stretch. Authors submitting work to the IEEE should note
% that the IEEE rarely uses double column equations and that authors should try
% to avoid such use. Do not be tempted to use the cuted.sty or midfloat.sty
% packages (also by Sigitas Tolusis) as the IEEE does not format its papers in
% such ways.
% Do not attempt to use stfloats with fixltx2e as they are incompatible.
% Instead, use Morten Hogholm'a dblfloatfix which combines the features
% of both fixltx2e and stfloats:
%
% \usepackage{dblfloatfix}
% The latest version can be found at:
% http://www.ctan.org/pkg/dblfloatfix

\usepackage{booktabs}


%\ifCLASSOPTIONcaptionsoff
%  \usepackage[nomarkers]{endfloat}
% \let\MYoriglatexcaption\caption
% \renewcommand{\caption}[2][\relax]{\MYoriglatexcaption[#2]{#2}}
%\fi
% endfloat.sty was written by James Darrell McCauley, Jeff Goldberg and 
% Axel Sommerfeldt. This package may be useful when used in conjunction with 
% IEEEtran.cls'  captionsoff option. Some IEEE journals/societies require that
% submissions have lists of figures/tables at the end of the paper and that
% figures/tables without any captions are placed on a page by themselves at
% the end of the document. If needed, the draftcls IEEEtran class option or
% \CLASSINPUTbaselinestretch interface can be used to increase the line
% spacing as well. Be sure and use the nomarkers option of endfloat to
% prevent endfloat from "marking" where the figures would have been placed
% in the text. The two hack lines of code above are a slight modification of
% that suggested by in the endfloat docs (section 8.4.1) to ensure that
% the full captions always appear in the list of figures/tables - even if
% the user used the short optional argument of \caption[]{}.
% IEEE papers do not typically make use of \caption[]'s optional argument,
% so this should not be an issue. A similar trick can be used to disable
% captions of packages such as subfig.sty that lack options to turn off
% the subcaptions:
% For subfig.sty:
% \let\MYorigsubfloat\subfloat
% \renewcommand{\subfloat}[2][\relax]{\MYorigsubfloat[]{#2}}
% However, the above trick will not work if both optional arguments of
% the \subfloat command are used. Furthermore, there needs to be a
% description of each subfigure *somewhere* and endfloat does not add
% subfigure captions to its list of figures. Thus, the best approach is to
% avoid the use of subfigure captions (many IEEE journals avoid them anyway)
% and instead reference/explain all the subfigures within the main caption.
% The latest version of endfloat.sty and its documentation can obtained at:
% http://www.ctan.org/pkg/endfloat
%
% The IEEEtran \ifCLASSOPTIONcaptionsoff conditional can also be used
% later in the document, say, to conditionally put the References on a 
% page by themselves.




% *** PDF, URL AND HYPERLINK PACKAGES ***
%
%\usepackage{url}
% url.sty was written by Donald Arseneau. It provides better support for
% handling and breaking URLs. url.sty is already installed on most LaTeX
% systems. The latest version and documentation can be obtained at:
% http://www.ctan.org/pkg/url
% Basically, \url{my_url_here}.




% *** Do not adjust lengths that control margins, column widths, etc. ***
% *** Do not use packages that alter fonts (such as pslatex).         ***
% There should be no need to do such things with IEEEtran.cls V1.6 and later.
% (Unless specifically asked to do so by the journal or conference you plan
% to submit to, of course. )

\usepackage[
backend=biber,
style=alphabetic,
sorting=ynt
]{biblatex}

\addbibresource{references.bib}


% correct bad hyphenation here
\hyphenation{op-tical net-works semi-conduc-tor}

\begin{document}

\title{Numerical Solutions to the Schrödinger Equation for the Hydrogen Molecule} 

\author{Louis Liu}% <-this % stops a space

\markboth{Louis Liu}%
{Shell \MakeLowercase{\textit{et al.}}:}

\maketitle

\begin{abstract}
    This paper describes the Variational Quantum Monte Carlo (VMC) method for generating numerical solutions to the time-independent Schrödinger equation for the Hydrogen molecule, $H_2$. A parametrised anstatz was optimised using .... [minimisation algorithm] .... to find the variational minima of the Hamiltonian expectation value for a range of interatomic distances, $r_0$ between the two nuclei. The produced energy profile for $H_2$ at different distances $r_0$ was fitted to a Morse potential, yielding estimates for the equilibrium bond length and dissociation energy of the molecule. The results produced unequivocally confirm the theoretical behaviour of the system, especially as the distance $r_0$ between the nuclei tends to zero, and the energy increases rapidly. The method and results described here present support for the scalability of stochastic simulation methods in atomic physics.
\end{abstract}


\section{Introduction}
    \IEEEPARstart{I}{\MakeLowercase{n}} many physical systems, the Schrödinger equation cannot be solved analytically, requiring the use of numerical approximation. Furthermore, in many multi-particle systems, the dimensionality of the Hilbert space scales exponentially with the number of particles. Variational Quantum Monte Carlo (VMC) is a useful technique for extracting energy predictions from the Schrödinger Equation in these cases by stochastic means. In many cases, it is most useful to to determine the ground state and its associated energy, which is obtained, fortunately, by minimising the expectation value of the Hamiltonian operator, $\braket{H}$.

    This paper will first consider the cases of the quantum harmonic oscillator, followed by the hydrogen atom, both of which have analytic solutions suitable for validation of the method. Finally, the case of the Hydrogen Molecule is considered, whose treatment demands further computational analysis.


    %This is then used to find the expected energy, $\braket{H}$, of the coordinates, which leads to a determination of the ground state energy of the system
\section{Theory} \label{Theory}
    
    
\section{Numerical Derivatives} 
    Finite difference methods were tested to obtain numerical estimates of the derivatives of the functions above. The second-order methods for computing the local energy were testsed first in one dimension using the simple harmonic oscillator, [see FIGURE]. The central difference method outperformed the backward and forward difference methods, with errors $\approx$3x smaller, whilst sharing the same optimal stepsize. Thus, the central difference method was adopted as the primary differentiator.

    Tests for single derivatives and in 3 dimensions for the H$_2$ molecule yielded similar results [FIGURE]. In the tested cases, it was clear that the optimal method was to 8th order, having the lowest error for the largest stepsize, with very little runtime difference even for wavefunctions. In the particular case of determining the ground state energies in the H$_2$ system, the Jastrow factor present in the anstatz made minimising the stepsize the primary goal when singularly differentiating the function. This was verified looking at the self-consistency of the function, which indicated that the second-order differentiator was most suitable.
    
\section{Stochastic Function Sampling}
    The methodology used to solve the systems aforementioned relies on the use of function sampling to calculate the value of the wavefunction at a given coordinate in space. Given their high dimensionality and nonlinearity, the wavefunctions we consider cannot be directly sample; thus, they are well-suited to rejection-based Markov Chain Monte Carlo (MCMC) type sampling methods.
    
    Here, we discuss the methods considered and tested for this implementation. 
    
    \subsection{Metropolis Hastings Sampling}
        One of the primary methods used in this simulation was the Metropolis Hastings (MH) sampling algorithm. This algorithm is particularly useful given use of the Jastrow-Slater ansatz in the H$_2$ scenario, as the acceptance relies only on the ratio of probability densities. Therefore, the wavefunction can be sampled without the need to calculate the non-trivial normalisation constant.

        In order to ensure effective sampling of the wavefunction, the proposal function was tuned to maintain the canonical acceptance rate of 0.234 [ref] [FIG? - include small justification on proposal function] whilst maintaining low autocorrelation. This was also confirmed visually and using STATISTICAL TESTS OF TARGET? initially for the quantum harmonic oscillator. 

        The number of steps in the sampler, $N_S$, is the main factor affecting both the numerical error and runtime of the simulation. Therefore its tuning is an important consideration. From visual analysis, sample sizes below 10,000 produced distributions with slight deviation from the true function, whilst sample sizes in excess of 100,000 represented it well. However, the difference in numerical error with the different sample numbers was evident. When computing the ground state energy for the Hydrogen atom, $N_S = 10,000$ produced an error of [ERROR], whilst an $N_S$ of $100,000$ produced only [ERROR]. Therefore, $N_S = 100,000$ was used in the optimisation of the hydrogen molecule due to the necessity of accurate numerical calculations. Due to its significance, further rigorous optimisation of this parameter is warranted. 

        The algorithm was initialised at 1 in each coordinate, since for all functions considered, the probability density at this point was non-zero. Due to the rapid decay of the wavefunction away from $r=0$ in all cases, a box of $\pm10$a.u. was used around the origin.

        The use of thinning substantially reduced autocorrelation, at the expense of fewer samples for a given runtime. [FIGURE]

    \subsection{Metropolis-adjusted Langevin algorithm}
        Improving upon the efficiency of the standard Metropolis-Hastings algorithm, the Metropolis-adjusted Langevin Algorithm (MALA) was implemented. Standard Metropolis sampling uses a diffusive "random walk" that scales poorly with system size, whilst MALA utilizes the gradient of the wavefunction to propose moves towards regions of higher probability. 

        Though this method undersamples low-probability regions, its efficacy is maintained for calculating the local energy by preferentially sampling regions that contribute more significantly. Additionally, for particular timesteps, MALA outperforms MH by acceptance rate and autocorrelation (without thinning). [FIGURE]

        Mode locking can occur when this method is used, where the sampler becomes trapped in a single peak if the peaks are spatially separated, as is of particular concern in the H$_2$ case. Adapting this method was therefore necessary, that we call the Stochastic Metropolis-adjusted Langevin algorithm (sMALA), in which a proportion of proposals are generated via a random walk. In this implementation, the proportion is limited to 0.1 to mitigate mode-locking.
        
    \subsection{Other Methods}   
        There are also some non-trivial examples of unsuitable methods.
        
        An example of a conceptually viable algorithm choice would be rejection sampling and its associated variants. However, the need to propose an efficient overlaying proposal distribution makes it hard to implement a general algorithm whilst maintaining computational efficiency. For a 6-dimensional, two-electron system, finding a bounding function that hugs the electron density tightly is tricky. Furthermore, as the dimensionality grows, the acceptance probability falls as the volume of the proposal distribution grows exponentially; hence, the usefulness of this method is limited to low dimensions.     

        Additionally, whilst the Gibbs sampling method facilitates a 100\% acceptance rate, it is not viable for this case. It requires a conditional probability distribution that is precluded by the non-linear Jastrow factor in the anstatz that we use for the hydrogen molecule.
        
        
\section{Computational Method}

\section{Experimental Results and Discussion}
    
\section{Conclusion}



% that's all folks
\end{document}


